\section{Discussion}
In this section, we discuss
open questions regarding the design of multi-agent systems for complex-tasks (Sec.~\ref{sec:The Multi-Agent Paradigm}),
current limitations (Sec.~\ref{sec:Limitations}),
and risks and risk mitigation for agents that autonomously operate computers (Sec.~\ref{sec:Risks and Mitigations}).

%%%%%%%%%%%%%%%%%%%%%%%%%%%%%%%%%%%%%%%%%%%%%%
\subsection{The Multi-Agent Paradigm}
\label{sec:The Multi-Agent Paradigm}

At the core of \team is its multi-agent design.
We believe that this design is a principal contributing factor to \team's performance.
Indeed, we observe that most other top-performing systems also follow a multi-agent design (Sec.~\ref{sec:Results}).

We argue that, beyond performance, the multi-agent setup offers numerous other advantages over the single-agent setup, in terms of ease-of-development, cost, and raw performance. 
%These benefits arise from the following properties:
For example, organizing skills into distinct agents can simplify development, much like with object-oriented programming. The separation of concerns across agents allows developers to focus model choices, prompting strategies, and other parameters to align to specific tasks (e.g., the web surfer agent benefits from multi-modality and structured output, but need not worry about writing code). Similarly, agent modularity can increase agent re-usability and ease of extensibility, particularly when teams are carefully designed to enable a plug-and-play approach. For example, \team's design facilitates adapting the team's functional scope by simply adding or removing agents, without requiring modifications to other agents's prompts or the overall flows and orchestration strategy. In contrast, monolithic single-agent systems often rely on constrained workflows that can be difficult to adapt or extend.

As a consequence of such modularity, each agent can be implemented in a fashion best suited for its purpose. In this paper, we leveraged this diversity to incorporate the o1-preview model into some roles (e.g., the Coder, and the outer loop  of the Orchestrator), while relying on a general purpose multi-modal model (GPT-4o) for web and file surfing. Looking ahead, we see this approach being used to reduce reliance on large models --  whereas some subtasks might require the largest language models available, others (e.g., grounding actions in WebSurfer, or summarizing large files in FileSurfer) might be amenable to much smaller -- and thus cheaper -- models. Different subtasks may also require different modalities, and some subtasks might be offloaded to traditional, non-AI, tools (e.g., code execution, for which a standard code execution environment is both sufficient and necessary). By embracing this diversity, multi-agent systems can become more performant at lower costs.


Understanding and quantifying the empirical advantages of multi- vs.\ single-agent setups constitutes a key question for future research.
Moreover, many variants of the multi-agent setup are possible. Here we opted for a single, centralized control flow pattern, where the Orchestrator agent plans for, and invokes, specialized worker agents.
Many other patterns are conceivable.
For instance, we might consider less centralized control flows, such as a peer-to-peer setting where each agent decides on its own which other agent should take control next.
At the other end of the spectrum, we might consider an even more rigid control flow where the orchestrator follows its own plan strictly (e.g., by encoding it as an executable program), rather than simply maintaining the plan it its prompt for chain-of-thought prompting.
Determining which control flow works best for which tasks is of considerable theoretical and practical importance.
%

In addition to the above-mentioned control-flow considerations, an alternate design dimension relates to the axes along which work is divided. \team's design diverges from other recent examples of multi-agent systems in that agents take on functional or tool-based responsibilities (web browser, computer terminal, etc.), rather than role-based responsibilities analogous to human teams (planner, researcher, data analyst, critic, etc.). In our experience, tool-centric agents can provide a cleaner separation of concerns compared to role-based agents, and a cleaner path to re-usability and compositionality -- if a web browser is a generic, multi-purpose tool, then a capable WebSurfer agent may hope to be generic and multi-purpose as well. Conversely, role-based patterns may require multiple agents to have redundant capabilities, while each agent fills only highly-specialized niche roles. For example, both a researcher and data analyst agent may need to operate a web browser or write code to complete their assigned tasks. Future work should empirically compare the performance of teams built with function- and role-based agents and examine the impact of each approach on the ease of development and debugging.





%%%%%%%%%%%%%%%%%%%%%%%%%%%%%%%%%%%%%%%%%%%%%%
\subsection{Limitations}
\label{sec:Limitations}
%%%%%%%%%%%%%%%%%%%%%%%%%%%%%%%%%%%%%%%%%%%%%%

Our work necessarily comes with certain limitations, some of which affect today's state of the field in general, and some of which are specific to our solution:



\begin{itemize}
    \item \textbf{Accuracy-focused evaluation:} Similar to other state-of-the-art systems, \team{} was evaluated on benchmarks that consider only the accuracy or correctness of final results. While considerably easier and more convenient to measure, such evaluations overlook  important considerations such as cost, latency, user preference and user value \cite{kapoor2024aiagentsmatter}. For example, even a partially correct trajectory may be valuable \cite{dibia-etal-2023-aligning}, whereas a perfectly accurate answer, delivered too late or at too high cost, may have no, or even negative, value. Designing evaluation protocols that incorporate these considerations, and that include subjective or open-ended tasks where correctness is less clear, remains an ongoing open-challenge in this the field.  
    \item \textbf{High cost and latency:} Although it was not part of the formal evaluation of \team, we would be remiss to skip mention of cost and latency in any discussion of limitations \cite{kapoor2024aiagentsmatter}. \team requires dozens of iterations and LLM calls to solve most problems. The latency and cost of those calls can be prohibitive, incurring perhaps several US dollars, and tens of minutes per task. We believe we can reduce these costs through targeted application of smaller local models, for example to support tool use in FileSurfer and WebSurfer, or set-of-mark action grounding in WebSurfer. Adding human oversight and humans-in-the-loop, may also save costs by reducing the number of iterations incurred when agents are stuck and problem-solving. This remains an active and ongoing area of future research.
    \item \textbf{Limited modalities:} \team cannot currently process or navigate all modalities. For example, WebSurfer cannot watch online videos -- though it often compensates by consulting transcripts or captions. Likewise, FileSurfer operates by converting all documents to Markdown, making it impossible to answer questions about a document's figures, visual presentation style, or layout. Audio files are similarly processed through a speech transcription model, so no agents can answer questions about music, or non-speech content. Benchmarks like GAIA exercise each of these skills. We would expect both benchmark and general task performance to improve with expanded support of multi-modal content. Future options include expanding \team's WebSurfer and FileSurfer agents' multi-modal capabilities or adding Audio and VideoSurfer agents specialized in handling audio and video processing tasks to the \team team. The latter approach is most inline with the value proposition of the multi-agent paradigm around easing development and reuse.      
    \item \textbf{Limited action space:} While agents in \team{} are afforded tools for the most common actions, tooling is not comprehensive. This simplifies the task of action grounding, but can lead to paths that are impossible to execute. For instance, the WebSurfer agent cannot hover over items on a webpage, or drag and resize elements. This can be limiting when interacting with maps, for example. Likewise, FileSurfer cannot handle all document types, and the Coder and Computer Terminal agents cannot execute code that requires API keys, or access to external databased or computational resources. We expect this limitation to close, over time, from two directions: first, we expect tool use to standardize across the industry, greatly enriching the set of tools available to agents. Second, similar to WebSurfer, agents will become better able to use operating systems and applications, affording them access to a broad range of tools developed for people.
    \item \textbf{Limited coding capabilities:} The \team{} Coder agent is particularly simple: it writes a new standalone Python program in respose to each coding request. In cases, where a prior coding attempt requires debugging, the Coder must correct the code by outputting an entirely new code listing. This is clearly not ideal. Two important limitations arise from this design choice: first, the Coder is ill-suited to operate over existing complex, or multi-file, code bases. Overcoming this limitation will be necessary to be competitive on benchmarks like SWE-bench \cite{jimenez2024swebenchlanguagemodelsresolve}. Second, the Coder sometimes fails because it expects functions that it previously defined to be available later in the workflow. Migrating to a Jupyter Notebook-like design, where later invocations simply add cells to a notebook, might mitigate this particular issue. This capability is presently supported by the AutoGen library upon which \team is built, and should be explored further.
    \item \textbf{Fixed team membership:} Additionally, \team's composition is fixed to a common set of five agents: Orchestrator, WebSurfer, FileSurfer, Coder, and ComputerTerminal. When agents are not needed, they simply serve as a distraction to the Orchestrator, and this may lower performance. Conversely, when extra expertise might be needed, it is simply unavailable. We can easily imagine an alternative approach where agents are added or removed dynamically, based on the task need. 
    \item \textbf{Limited learning:} Finally, although \team can adapt its strategy based on trial and error within a single task attempt, such insights are discarded between tasks. We observe the direct consequences of this design in WebArena, where many problems share a common set of core sub-tasks (e.g., finding a particular thread or user profile). When competing on this benchmark, \team's agents need to discover and rediscover solutions to these sub-tasks over and over. This is exhausting and frustrating to watch, is highly prone to error, and can incur significant additional costs. Overcoming this limitation through long-term memory is a key direction for future research.
\end{itemize}



%%%%%%%%%%%%%%%%%%%%%%%%%%%%%%%%%%%%%%%%%%%%%%
\subsection{Risks and Mitigations}
\label{sec:Risks and Mitigations}
%%%%%%%%%%%%%%%%%%%%%%%%%%%%%%%%%%%%%%%%%%%%%%

The agents described in this paper interact with a digital world designed for, and inhabited by, humans. This carries inherent and undeniable risks. In our work we mitigate such risks by running all tasks in containers, leveraging synthetic environments like WebArena, choosing models with strong alignment and pre- and post-generation filtering, and by closely monitoring logs during and after execution. Nevertheless, we observed the agents attempt steps that would otherwise be risky. For example, during development, a mis-configuration prevented agents from successfully logging in to a particular WebArena website. The agents attempted to log in to that website until the repeated attempts caused the account to be temporarily suspended. The agents then attempted to reset the account's password. In other cases, agents recognized that the WebArena Postmill website was not Reddit, then directed agents to the real website to commence work -- this was ultimately blocked by network-layer restrictions we had put in place. Likewise, we observed cases where agents quickly accepted cookie agreements and website terms and conditions without any human involvement (though captchas were correctly refused). More worryingly, in a handful of cases -- and until prompted otherwise -- the agents occasionally attempted to recruit other humans for help (e.g., by posting to social media, emailing textbook authors, or, in one case, drafting a freedom of information request to a government entity). In each of these cases, the agents failed because they did not have access to the requisite tools or accounts, and/or were stopped by human observers. To this end, it is imperative that agents operate under a strict principle of least privilege, and maximum oversight.

In addition to these observed and mitigated risks, we can anticipate new risks from such agentic systems on the horizon. As an example, as agents operate on the public internet, it is possible that they are subject to the same phishing, social engineering, and misinformation attacks that target human web surfers. The fact that the WebSurfer only occasionally recognized that Postmill was not Reddit, in WebArena, lends credence to the concern that agents can be fooled. If agents are equipped with a user's personal information  -- for example, to complete tasks on their behalf --  then this could potentially put that information at risk. Moreover, we can imagine that such attacks may be made more reliable and effective if attackers anticipate agentic use, and seed external material with specially crafted instructions or prompt-injections. To address these challenges, we can imagine several mitigations such as increasing human oversight, equipping agents with tools to validate external information (e.g., checking for typo-squatting in URLs, and ensuring TLS certificates, etc.), and including phishing rejection examples and other web-savvy skills in a model's post-training and instruction tuning. 

Another cross-cutting mitigation we anticipate becoming important is equipping agents with an understanding of which actions are easily reversible, which are reversible with some effort, and which cannot be undone. As an example, deleting files, sending emails, and filing forms, are unlikely to be easily reversed. This concept is explored in some detail in \cite{zhang2024interaction}, which provides a compelling framework for considering agent safety. When faced with a high-cost or irreversible action, systems should be designed to pause, and to seek human input.

Recent research has also investigated how interactions between multiple agents, such as iterative requests, long contexts, or cascading errors, may impact the effectiveness of the existing model alignment and guardrails upon which we rely. For example, the crescendo multi-turn attack \cite{russinovich2024great}, operates by prompting the model with a benign request, the slowly escalates the requests, building a pattern of agent compliance, until finally asking for a response that would otherwise be refused. In a multi-agent system, it is possible that a malicious agent, or an unfortunate accident, could result in a similar pattern of escalation. Fortunately, strong pre- and post-filtering, on the prompts and responses, remain a reasonable mitigation to such risks for the short-term. Moving forward, we strongly encourage model alignment work to focus on multi-turn scenarios. We also believe that red-team exercises are imperative to identify and mitigate such risks. 

Finally, there are potential long-term societal impacts of agentic systems, such as the potential to deskill or replace workers, leading to potential economic disruption. We believe it is therefore critical to work towards designing systems that facilitate effective collaboration between people and agents, such that humans and agents working together can achieve more than agents working alone. 
